\documentclass[11pt]{article}

\usepackage[T1]{fontenc}
\usepackage{lmodern}
\usepackage[margin=1in]{geometry}
\usepackage{amsmath,amssymb,amsthm,mathtools,bm}
\usepackage{booktabs}
\usepackage{array}
\usepackage{enumitem}
\usepackage{float}
\usepackage{xcolor}
\usepackage{microtype}
\usepackage{hyperref}
\usepackage[nameinlink,capitalise,noabbrev]{cleveref}

\hypersetup{
  colorlinks=true,
  linkcolor=blue!55!black,
  citecolor=blue!55!black,
  urlcolor=blue!55!black,
  pdftitle={Operators as Boundary Spectral Relations},
  pdfauthor={Andrew Kelleher}
}

\allowdisplaybreaks
\numberwithin{equation}{section}

\newtheorem{theorem}{Theorem}[section]
\newtheorem{proposition}[theorem]{Proposition}
\newtheorem{corollary}[theorem]{Corollary}
\newtheorem{lemma}[theorem]{Lemma}
\theoremstyle{definition}
\newtheorem{definition}[theorem]{Definition}
\newtheorem{remark}[theorem]{Remark}
\theoremstyle{remark}
\newtheorem{openproblem}[theorem]{Open problem}

\newcommand{\C}{\mathbb{C}}
\newcommand{\R}{\mathbb{R}}
\newcommand{\id}{\mathrm{id}}
\newcommand{\spec}{\operatorname{spec}}
\newcommand{\ran}{\operatorname{ran}}
\newcommand{\tr}{\operatorname{tr}}
\newcommand{\diag}{\operatorname{diag}}
\newcommand{\dist}{\operatorname{dist}}
\newcommand{\geo}{\operatorname{geo}}
\newcommand{\Span}{\operatorname{span}}
\newcommand{\ket}[1]{\lvert #1\rangle}
\newcommand{\bra}[1]{\langle #1\rvert}
\newcommand{\braket}[2]{\langle #1\mid #2\rangle}
\newcommand{\matrixel}[3]{\langle #1\mid #2\mid #3\rangle}
\newcommand{\norm}[1]{\left\lVert #1\right\rVert}

\title{Operators as Boundary Spectral Relations:\\
Exact Magnetic-Graph Synthesis and Fixed-Boundary Cobordism Relaxation}
\author{Andrew Kelleher\\Twin Vector Labs LLC}
\date{Draft of August 26, 2026}

\begin{document}

\maketitle

\begin{abstract}
We study a spectral encoding of a finite-dimensional linear operator in which the
operator is not identified with one fitted eigenvector.  Instead, the graph of a
map $U:\mathcal H_-\to\mathcal H_+$ is required to equal the boundary trace of a
common eigenspace of a weighted magnetic connection Laplacian.  This formulation
makes input substitution a consequence of linearity: fitting a basis at one
common eigenvalue fits every superposition, provided that the boundary trace has
no additional modes and its input projection is nonsingular.  We prove an exact
finite-network realization theorem.  For any Hermitian boundary matrix, any
fixed induced boundary magnetic graph, and any positive spectral parameter, a
finite simple magnetic graph can be constructed whose Schur response is exactly
the prescribed matrix while every intrinsic boundary edge remains unchanged.
Applying the theorem to a canonical graph-of-map penalty realizes every complex
linear map, including singular and non-normal maps.  Charge conservation is
characterized separately as equivariance of the boundary relation; it is a
symmetry condition, not an existence condition.  Boundary permutations are
shown to be framing data and therefore part of operator identification.

We then describe a fixed-boundary relaxation implemented in Tessera.  The method
pins incoming and outgoing restrictions for a complete training basis, varies
only bulk geometry and bulk witness amplitudes, and minimizes the full-complex
connection-Laplacian eigenresidual at a common eigenvalue.  Deterministic tests
recover a generic two-dimensional unitary to relative error
$2.51\times10^{-16}$ with zero recorded boundary drift.  A completed 77-case
stress test verifies 43 cases, exhausts the optimization budget in 22, rejects
eight preparations under the current normalization and isolated-eigenstate
preconditions, and exposes four identifiability or
common-eigenvalue failures.  These are numerical solver outcomes, not a
universality proof.  The exact theorem is presently a magnetic-graph theorem.
Proving that the construction can always be thickened into a nondegenerate
triangulated cobordism without changing its boundary response remains the
central geometric problem.  Accordingly, the work establishes an exact
operator-as-spectral-relation framework and a constructive graph realization,
but does not claim a topological quantum field theory or a universal simplicial
realization theorem.
\end{abstract}

\noindent\textbf{Keywords.}
inverse spectral problem; magnetic graph Laplacian; Schur complement;
cobordism; quantum operator; charge equivariance; fixed-boundary relaxation

\section{Introduction}

A cobordism $W:M_-\rightsquigarrow M_+$ is geometrically interpreted as a
process between an incoming and an outgoing boundary.  In a topological quantum
field theory (TQFT), a functor assigns a linear map to $W$, and gluing
cobordisms corresponds to composition of maps~\cite{atiyah1988}.  The inverse
question is less standard: can a prescribed linear map be represented by
finite geometry, and can that representation be found by relaxation while the
boundary data remain fixed?

The motivating numerical experiments use weighted simplicial complexes.  A
boundary state is represented in a selected spectral frame on a boundary
component.  Given an algebraic operator $U$, the expected outgoing state is
computed first.  Incoming and outgoing restrictions are then pinned, while
bulk weights, connection phases, topology, and auxiliary interior amplitudes
are adjusted to reduce a full-complex eigenresidual.  Once one geometry has
been fitted using a complete basis, held-out superpositions are evaluated with
the same geometry.

There are three materially different claims that must not be conflated.

\begin{enumerate}[label=(\roman*)]
  \item A single vector proportional to $\operatorname{vec}(U)$ can be prepared
  as an eigenstate.  This is state synthesis in a Choi-like coordinate system,
  not by itself an operational representation of $U$.
  \item Boundary traces of one common eigenspace can form the graph
  $\{(x,Ux):x\in\mathcal H_-\}$.  This is an operational representation:
  changing $x$ changes the outgoing restriction to $Ux$ by linearity.
  \item A scalar weighted magnetic network realizing that spectral relation can
  always be chosen to be the one-skeleton of a valid triangulated cobordism with
  the prescribed frozen boundary.  This additional relative-realization claim
  is open.
\end{enumerate}

This paper proves the second claim abstractly and gives an exact finite magnetic
graph construction.  It reports numerical evidence for the corresponding
fixed-boundary relaxation, while isolating the unproved step needed for the
strongest simplicial-cobordism interpretation.

\subsection{Contributions}

The principal contributions are as follows.

\begin{enumerate}[leftmargin=2.2em]
  \item We define an operator represented by a cobordism as a \emph{boundary
  spectral relation}, rather than as a distinguished bulk vector.
  \item We prove that complete-basis witnesses at a common eigenvalue extend to
  every superposition, and state the rank condition required to exclude
  spurious boundary modes.
  \item We give a canonical positive semidefinite matrix $K_U$ whose kernel is
  exactly the graph of an arbitrary complex linear map $U$.
  \item We prove constructively that every Hermitian matrix is the boundary
  Schur response of a finite simple scalar magnetic graph at any prescribed
  positive spectral parameter, without changing the induced boundary graph.
  \item We prove that charge conservation is equivalent to invariance of the
  graph relation, and to commutation of $K_U$ with total boundary charge.
  \item We identify attachment permutations as boundary framings and give
  residual-, gap-, and conditioning-based certification criteria.
  \item We document the exact semantics and current numerical limits of the
  Tessera relaxation experiments.
\end{enumerate}

\section{Relation to prior work}

The functorial assignment of state spaces and linear maps to manifolds and
cobordisms is axiomatized by Atiyah~\cite{atiyah1988}.  Our direction is inverse:
the map is specified first, and a finite spectral network is synthesized to
carry its graph.  We do not establish the identity, gluing, monoidal, or
diffeomorphism-invariance axioms required of a TQFT.

Encoding an operator as a vector is familiar from the
Jamio\l kowski--Choi correspondence~\cite{jamiolkowski1972,choi1975}.  That
correspondence is useful for preparing training targets, but a fitted vector is
not automatically a device that accepts a new input.  The graph-of-map
formulation used here supplies the missing input/output relation.

Eliminating interior vertices by a Schur complement is Kron reduction
\cite{dorfler2013}.  Closely related inverse response questions have been
studied for metric quantum graphs: Friedlander showed that arbitrary real
symmetric matrices can occur as Dirichlet-to-Neumann maps for a quantum-graph
Schr\"odinger operator~\cite{friedlander2017}.  Inverse eigenvalue questions for
ordinary graph Laplacians remain constrained by graph structure
\cite{fallat2024}.  Magnetic phases enlarge the realizable class from real
symmetric to complex Hermitian matrices; magnetic and connection Laplacians
are standard objects in spectral graph theory and data analysis
\cite{fabila2018,singer2012}.  Our finite construction is explicit, preserves a
preassigned induced boundary graph, and is tailored to a graph-of-map kernel.

The triangulated extension is related to, but not resolved by, bistellar
equivalence of triangulations~\cite{pachner1991}.  Bistellar or stellar moves
control topology; they do not automatically preserve a prescribed magnetic
Schur response.

\section{The discrete spectral object}
\label{sec:operator}

\subsection{Scalar magnetic connection Laplacian}

Let $G=(V,E)$ be a finite simple graph.  Give each unoriented edge
$\{u,v\}$ a weight $w_{uv}>0$ and an oriented phase
$a_{uv}=e^{i\theta_{uv}}\in U(1)$, with $a_{vu}=\overline{a_{uv}}$.  The scalar
magnetic connection Laplacian $L_G$ acts on $z\in\C^V$ by
\begin{equation}
  (L_Gz)_u=\sum_{v:\{u,v\}\in E}w_{uv}(z_u-a_{uv}z_v).
  \label{eq:magnetic-laplacian}
\end{equation}
It is Hermitian and positive semidefinite because
\begin{equation}
  z^\dagger L_Gz
  =\sum_{\{u,v\}\in E}w_{uv}\lvert z_u-a_{uv}z_v\rvert^2.
  \label{eq:energy}
\end{equation}
In the Tessera implementation, $w_{uv}$ is the positive quantity derived from
the squared edge length and $\theta_{uv}$ is the optimized real connection
phase.

This distinction is important: the experiments analyzed here use the
degree-zero $\C^*$ connection Laplacian in~\eqref{eq:magnetic-laplacian}.  They
do not use a degree-one Hodge Laplacian, and they do not minimize a Regge action
or a harmonic-period residual.  The ordinary degree-zero discrete exterior
calculus Laplacian is a separate operator in the codebase.

\subsection{Boundary response}

Choose a nonempty boundary vertex set $B\subseteq V$ and write
$I=V\setminus B$.  Relative to $\C^V=\C^B\oplus\C^I$, write
\begin{equation}
  L_G=\begin{pmatrix}A&C\\ C^\dagger&D\end{pmatrix}.
  \label{eq:block-laplacian}
\end{equation}
For a real spectral parameter $\lambda\notin\spec(D)$, define the boundary
Schur response
\begin{equation}
  S_G(\lambda)
  =A-\lambda I_B-C(D-\lambda I_I)^{-1}C^\dagger.
  \label{eq:schur}
\end{equation}

\begin{definition}[Boundary spectral relation]
The boundary spectral relation of $(G,B)$ at $\lambda$ is
\begin{equation}
  \mathcal R_G(\lambda)
  =\{z|_B:z\in\ker(L_G-\lambda I_V)\}\subseteq\C^B.
  \label{eq:relation}
\end{equation}
\end{definition}

\begin{proposition}[Schur characterization]
\label{prop:schur}
If $\lambda\notin\spec(D)$, then
\begin{equation}
  \mathcal R_G(\lambda)=\ker S_G(\lambda),
  \qquad
  z_I=-(D-\lambda I_I)^{-1}C^\dagger z_B.
  \label{eq:schur-kernel}
\end{equation}
In particular, every admissible boundary trace has a unique interior lift.
\end{proposition}

\begin{proof}
The interior block of $(L_G-\lambda I)z=0$ gives the second identity.
Substitution into the boundary block gives
$S_G(\lambda)z_B=0$.  The converse follows by reversing these steps.
\end{proof}

\section{An operator is the graph of a boundary relation}

Let $\mathcal H_-\cong\C^d$ and $\mathcal H_+\cong\C^m$ be incoming and
outgoing state spaces, and let $U:\mathcal H_-\to\mathcal H_+$ be linear.  Its
graph is
\begin{equation}
  \Gamma_U=\{x\oplus Ux:x\in\mathcal H_-\}
  \subseteq\mathcal H_-\oplus\mathcal H_+.
  \label{eq:graph-U}
\end{equation}

Physical boundary vertices need not be coordinate vectors of these spaces.
Let $J_-:\mathcal H_-\hookrightarrow\C^{B_-}$ and
$J_+:\mathcal H_+\hookrightarrow\C^{B_+}$ be fixed injective state charts on
two disjoint boundary components.  For isometric charts, the embedded graph is
the range of
\begin{equation}
  T_Ux=J_-x\oplus J_+Ux.
  \label{eq:embedded-graph}
\end{equation}

\begin{definition}[Exact spectral realization]
Let $B=B_-\sqcup B_+$.  A framed magnetic graph $(G,B,J_-,J_+)$ realizes $U$
at $\lambda$ if
\begin{equation}
  \mathcal R_G(\lambda)=\ran T_U.
  \label{eq:exact-realization}
\end{equation}
Equivalently, after applying the inverse charts, the relation is $\Gamma_U$.
\end{definition}

The equality in~\eqref{eq:exact-realization}, rather than mere containment, is
the identifiability requirement.  It rules out an additional eigenvector with
zero incoming trace or a second outgoing trace for the same input.

\subsection{Why a complete basis generalizes}

\begin{theorem}[Basis-span transfer]
\label{thm:basis-span}
Let $b_1,\ldots,b_d$ be a basis of $\mathcal H_-$.  Suppose that for one graph
$G$ and one common eigenvalue $\lambda$, there are witnesses
$z_1,\ldots,z_d\in\ker(L_G-\lambda I)$ satisfying
\begin{equation}
  z_j|_{B_-}=J_-b_j,
  \qquad
  z_j|_{B_+}=J_+Ub_j.
  \label{eq:witness-restrictions}
\end{equation}
Then $\ran T_U\subseteq\mathcal R_G(\lambda)$.  Consequently every
$x\in\mathcal H_-$ has a full-graph eigenvector whose incoming and outgoing
restrictions are $J_-x$ and $J_+Ux$, respectively.  If additionally
$\dim\mathcal R_G(\lambda)=d$, then $G$ realizes $U$ exactly.
\end{theorem}

\begin{proof}
Write $x=\sum_j\alpha_jb_j$ and set $z=\sum_j\alpha_jz_j$.  Linearity gives
$(L_G-\lambda I)z=0$, and restriction gives
$z|_B=T_Ux$.  Thus $\ran T_U\subseteq\mathcal R_G(\lambda)$.  Since $T_U$ is
injective, $\dim\ran T_U=d$; equality follows under the stated dimension
condition.
\end{proof}

\begin{remark}[The eigenvector is global]
The output is not, in general, an eigenvector supported only on $B_+$.  It is
the outgoing restriction of a full-graph eigenvector.  The bulk amplitudes are
the unique lift in~\cref{prop:schur} when the interior block is nonsingular at
$\lambda$.
\end{remark}

\begin{remark}[A common eigenvalue is essential]
If witness $z_j$ belongs to eigenvalue $\lambda_j$ and the $\lambda_j$ are not
equal, a generic linear combination is not an eigenvector.  Small individual
residuals therefore do not establish operator transfer unless the eigenvalue
spread is also controlled.
\end{remark}

\subsection{Recovering the represented matrix}

Let $N\in\C^{(|B_-|+|B_+|)\times d}$ contain a basis of the boundary relation,
expressed in the selected charts, and split it as
\begin{equation}
  N=\begin{pmatrix}X\\Y\end{pmatrix}.
\end{equation}
The relation is the graph of a unique operator if and only if $X$ is square and
nonsingular; in that case
\begin{equation}
  U=YX^{-1}.
  \label{eq:operator-readout}
\end{equation}
This rank test is the correct operational test.  A vectorized bulk state is not
itself a map.  If desired, one may vectorize the recovered matrix in
\eqref{eq:operator-readout} to obtain its Jamio\l kowski--Choi vector or state.

\section{A universal algebraic target}

There is a canonical Hermitian matrix whose kernel is exactly $\Gamma_U$.  Set
\begin{equation}
  F_U=\begin{pmatrix}-U&I_m\end{pmatrix},
  \qquad
  K_U=F_U^\dagger F_U
  =\begin{pmatrix}
    U^\dagger U&-U^\dagger\\
    -U&I_m
  \end{pmatrix}.
  \label{eq:KU}
\end{equation}
For $z=x\oplus y$,
\begin{equation}
  z^\dagger K_Uz=\norm{y-Ux}^2,
  \qquad
  \ker K_U=\Gamma_U.
  \label{eq:KU-energy}
\end{equation}
The nonzero eigenvalues of $K_U$ are the eigenvalues of
$F_UF_U^\dagger=I_m+UU^\dagger$, so its nonzero spectral gap is at least one.
If $U$ is unitary, $K_U^2=2K_U$ and $K_U/2$ is the orthogonal projector onto
$\Gamma_U^\perp$.

For embedded isometric charts, a coordinate-free target on the full boundary
space is
\begin{equation}
  P_U=I_B-T_U(T_U^\dagger T_U)^{-1}T_U^\dagger,
  \qquad
  \ker P_U=\ran T_U.
  \label{eq:embedded-projector}
\end{equation}
Thus universal operator encoding reduces to realizing a prescribed Hermitian
boundary response with a prescribed kernel.

\section{Charge symmetry and boundary framing}

\subsection{Charge is equivariance, not existence}

Let $Q_-$ and $Q_+$ be Hermitian charge operators on $\mathcal H_-$ and
$\mathcal H_+$, and set $Q_\partial=Q_-\oplus Q_+$.  The usual
charge-conservation condition is
\begin{equation}
  Q_+U=UQ_-.
  \label{eq:charge-intertwiner}
\end{equation}

\begin{theorem}[Equivalent charge criteria]
\label{thm:charge}
For any linear $U$, the following statements are equivalent:
\begin{enumerate}[label=(\alph*)]
  \item $Q_+U=UQ_-$;
  \item $Q_\partial\Gamma_U\subseteq\Gamma_U$;
  \item $[K_U,Q_\partial]=0$.
\end{enumerate}
\end{theorem}

\begin{proof}
For $x\oplus Ux\in\Gamma_U$, invariance under $Q_\partial$ is precisely the
condition $Q_+Ux=UQ_-x$ for every $x$.  This proves (a)$\Leftrightarrow$(b).
The lower-left block of $[K_U,Q_\partial]$ vanishes exactly when
$Q_+U=UQ_-$.  Under this identity and its adjoint,
$[U^\dagger U,Q_-]=0$, so all remaining blocks vanish.  Hence
(a)$\Leftrightarrow$(c).
\end{proof}

The existence theorem below does not require~\eqref{eq:charge-intertwiner}.
Charge conservation becomes necessary only when the realizing geometry is also
required to respect the corresponding boundary symmetry.  In particular, a
charge-changing Pauli $X$ map is a valid spectral relation even though it is
not an intertwiner for the computational $Z$ charge.  Inside a degenerate
sector on which $Q_-=Q_+=qI$, every operator is automatically equivariant; the
charge condition then imposes no additional restriction.

\subsection{Attachment permutations are framing data}

The matrices assigned to a boundary depend on ordered state charts.  Under
invertible changes of incoming and outgoing frame $P_-$ and $P_+$, the same
abstract relation has coordinate matrix
\begin{equation}
  U\longmapsto P_+UP_-^{-1}.
  \label{eq:frame-transform}
\end{equation}
Vertex permutations are special cases.  Therefore the way an input complex is
attached to the bulk is not an incidental implementation detail: it specifies
the boundary framing in which the operator is read.

Trying every attachment permutation is a valid way to identify an unknown
framing.  Once selected, however, the same permutation must be held fixed for
the complete training basis and every held-out input.  Selecting a new
permutation separately for each state would fit a state-dependent rule rather
than one operator.  Moreover, a successful permutation does not by itself
resolve rank deficiency: the input block $X$ in
\eqref{eq:operator-readout} must remain well conditioned.

\section{Exact synthesis of arbitrary boundary response}
\label{sec:synthesis}

We now show that the algebraic target in~\cref{eq:KU} can always be realized by
a finite scalar magnetic network.  The result preserves every intrinsic edge
of a preassigned boundary graph.  New boundary-to-interior edges contribute to
the full diagonal block, as they must, but do not alter the induced boundary
subgraph.

\begin{theorem}[Fixed-boundary magnetic response synthesis]
\label{thm:response-synthesis}
Let $B$ be a finite nonempty vertex set, let $G_0$ be any weighted magnetic
graph induced on $B$, let $R\in\C^{B\times B}$ be Hermitian, and let $c>0$.
There exists a finite simple weighted magnetic graph $G$ with boundary $B$ such
that:
\begin{enumerate}[label=(\roman*)]
  \item the subgraph induced by $B$ is exactly $G_0$, with all of its weights
  and phases unchanged;
  \item $c\notin\spec(L_{II})$; and
  \item $S_G(c)=R$.
\end{enumerate}
If desired, $G$ may be chosen connected.
\end{theorem}

\begin{proof}
Let $L_0$ denote the magnetic Laplacian contribution of the fixed boundary
graph, embedded in the boundary block.  We construct mutually nonadjacent
interior vertices, so $L_{II}$ will be diagonal.

For each unordered pair $i<j$, define the required off-diagonal correction
\begin{equation}
  \delta_{ij}=R_{ij}-(L_0)_{ij}.
\end{equation}
If $\delta_{ij}\ne0$, add an interior vertex $v_{ij}$ connected only to $i$ and
$j$.  Give both edges a common weight $t\in(0,c/2)$ and choose their relative
phase to be $\arg\delta_{ij}$.  Since the diagonal entry of $L_{II}$ at this
vertex is $2t$, its off-diagonal contribution to the Schur response is
\begin{equation}
  \frac{t^2}{c-2t}e^{i\arg\delta_{ij}}.
  \label{eq:pair-response}
\end{equation}
For $a=\lvert\delta_{ij}\rvert$, the choice
\begin{equation}
  t=\sqrt{a^2+ca}-a
  \label{eq:pair-weight}
\end{equation}
lies in $(0,c/2)$ and satisfies $t^2/(c-2t)=a$.  Hence the pair gadget produces
exactly $\delta_{ij}$ and its conjugate in the two off-diagonal positions.  Its
contribution to each incident diagonal is
\begin{equation}
  g_c(t)=t+\frac{t^2}{c-2t}>0.
  \label{eq:pair-diagonal}
\end{equation}

After all pair gadgets have been inserted, only the diagonal entries may be
incorrect.  A leaf joined only to boundary vertex $i$ with edge weight
$p>0$, $p\ne c$, contributes
\begin{equation}
  f_c(p)=p-\frac{p^2}{p-c}=\frac{pc}{c-p}
  \label{eq:leaf-response}
\end{equation}
to the $i$th response diagonal.  The range of $f_c$ on $(0,c)$ is
$(0,\infty)$, and its range on $(c,\infty)$ is $(-\infty,-c)$.  Thus every
real diagonal correction $\eta$ is a sum of at most two leaf responses.  If
$\eta>0$ or $\eta<-c$, one leaf suffices, with
\begin{equation}
  p=\frac{c\eta}{c+\eta}>0.
\end{equation}
If $-c\leq\eta<0$, write $\eta=2c+(\eta-2c)$; the first term lies in
$(0,\infty)$ and the second in $(-\infty,-c)$, so two leaves suffice.  No leaf
is needed for $\eta=0$.

The interior diagonal entries are either $2t<c$ or $p\ne c$, so
$c\notin\spec(L_{II})$.  The construction uses distinct interior vertices and
no parallel edges, hence $G$ is simple, and no edge of $G_0$ was changed.

If the resulting graph is disconnected, choose a spanning tree on its boundary
components.  For every tree edge, insert two pair gadgets having the same
weight and opposite relative phases.  Their off-diagonal responses cancel,
while they connect the components.  Compensate their diagonal contributions
with the leaf construction above.  This preserves $S_G(c)=R$ and yields a
connected graph.
\end{proof}

\begin{corollary}[Universal graph realization of a linear map]
\label{cor:universal-map}
For every complex linear map $U:\C^d\to\C^m$, every fixed induced boundary
magnetic graph on $d+m$ coordinate vertices, and every $c>0$, there is a finite
simple magnetic graph whose boundary spectral relation at $c$ is exactly
$\Gamma_U$.  The same conclusion holds in fixed injective boundary charts by
using the target $P_U$ in~\cref{eq:embedded-projector}.
\end{corollary}

\begin{proof}
Apply~\cref{thm:response-synthesis} with $R=K_U$ in direct coordinates, or with
$R=P_U$ in embedded charts.  By~\cref{prop:schur}, the boundary relation is the
kernel of the response, which is the desired graph.
\end{proof}

The construction is an existence proof, not an efficient optimizer.  It uses
at most one degree-two interior vertex for each nonzero off-diagonal correction
and at most two leaves per boundary vertex, before optional connectivity
gadgets.  It applies to rectangular, singular, non-normal, charge-changing, and
poorly conditioned operators.

\section{What remains to obtain a simplicial-cobordism theorem}

Every finite graph is a one-dimensional simplicial complex.  That observation
does not prove the geometric claim of primary interest.  A one-dimensional
complex with branching interior vertices is generally not a cobordism between
the chosen boundary components, and thickening it into a higher-dimensional
triangulated manifold introduces additional edges and diagonal contributions
that may change~\eqref{eq:schur}.

\begin{openproblem}[Relative simplicial realization]
\label{open:relative-simplicial}
Given the magnetic graph produced by~\cref{thm:response-synthesis}, construct a
nondegenerate triangulated cobordism $W$ with the prescribed frozen boundary
such that the degree-zero connection Laplacian of $W$ has the same boundary
Schur response at $c$, possibly after adding response-neutral interior gadgets.
\end{openproblem}

A proof would upgrade~\cref{cor:universal-map} from finite discrete geometry to
the specific class of simplicial cobordisms used by the numerical machinery.
It must address four constraints simultaneously:

\begin{enumerate}[label=(\arabic*)]
  \item manifold links at every interior and boundary simplex;
  \item nondegenerate simplex geometry and compatible shared edge lengths;
  \item a fixed induced boundary triangulation and fixed boundary phases; and
  \item exact preservation, or controllably accurate approximation, of the
  Schur response at the selected spectral parameter.
\end{enumerate}

Pachner connectivity alone does not solve the problem, because a sequence of
bistellar moves need not preserve spectral response.  A possible alternative
for unitary maps is a rank-$d$ non-Abelian connection on a cylinder whose
holonomy is $U$.  That approach makes the operator geometrically immediate but
changes the scalar $U(1)$ machinery used in the present experiments; arbitrary
nonunitary maps would additionally require dilation or a nonunitary bundle
connection.

\section{Fixed-boundary relaxation}
\label{sec:relaxation}

The constructive proof above was derived to clarify what the numerical solver
is attempting to discover.  The present Tessera experiment uses two disjoint
triangulated boundary components and a shared bulk.  Its reference fixture
places a three-edge cycle on each boundary.  At unit weights and zero phases,
the cycle is $K_3$, whose connection-Laplacian spectrum is $\{0,3,3\}$.  The
two nonconstant Fourier modes provide a two-dimensional neutral state chart;
an optional constant mode has eigenvalue zero.

For a chosen basis $b_1,\ldots,b_d$, the expected outputs $Ub_j$ are computed
algebraically.  Their boundary representations are fixed as in
\cref{eq:witness-restrictions}.  The solver then varies bulk edge weights,
bulk connection phases, bulk witness amplitudes, and---when enabled---bulk
stellar growth.  Intrinsic boundary edge weights and phases remain bitwise
fixed, and every witness retains its exact incoming and outgoing restriction.

If $\widehat z_j$ denotes a full-complex witness and $\bar\lambda$ the common
Rayleigh target, the coupled objective is
\begin{equation}
  \mathcal J(W,\{\widehat z_j\})
  =\sum_{j=1}^d
  \norm{L_W\widehat z_j-\bar\lambda\widehat z_j}^2.
  \label{eq:relaxation-objective}
\end{equation}
Optimization uses bounded positive bulk weights, bounded real phases,
Levenberg--Marquardt restarts, and optional boundary-preserving stellar growth.
No Regge-stationarity term, curvature action, period objective, or explicit
charge penalty appears in~\eqref{eq:relaxation-objective}.

The protocol is therefore:

\begin{enumerate}[leftmargin=2.2em]
  \item prepare each boundary chart in isolation and record the intended
  degenerate spectral sector;
  \item select $U$ and compute $Ub_j$ algebraically for a complete basis;
  \item pin all incoming and outgoing restrictions and freeze intrinsic
  boundary geometry;
  \item relax only the shared bulk variables against the full-$W$ residual;
  \item certify common eigenvalue, boundary invariance, graph dimension, input
  rank, operator readout, and held-out superpositions; and
  \item when attachment framing is unknown, enumerate its permutations, but
  retain one selected framing for all training and test states.
\end{enumerate}

This is a supervised inverse spectral construction.  The target outputs are
used during fitting.  After fitting, basis-span transfer is exact in principle,
but only if the common-eigenvalue and identifiability hypotheses of
\cref{thm:basis-span} have been established.

\section{Numerical evidence}
\label{sec:numerics}

All values in this section are solver records, not mathematical assumptions.
The boundary-drift and restriction metrics are evaluated independently of the
optimization target.  Every eigenresidual field, including a held-out full-$W$
residual, is a squared norm.  Boundary and operator errors are ordinary norms
unless labeled otherwise.  Therefore a recorded objective $r$ certifies a
residual norm no larger than $\sqrt r$ for each witness.

\subsection{Historical direct-state mode}

The earlier experiment pinned one vectorized operator target.  It established
that the relaxation can synthesize both charge-preserving and charge-changing
states, but it did not establish input substitution.  Representative outcomes
are shown in~\cref{tab:historical}.

\begin{table}[htbp]
  \centering
  \caption{Historical direct-state synthesis.  Residual is the recorded
  squared objective.}
  \label{tab:historical}
  \begin{tabular}{lrrrr}
    \toprule
    Target & $\norm{[U,Q]}_F$ & residual & growth & boundary drift\\
    \midrule
    phase gate & $0$ & $7.43\times10^{-28}$ & 3 & $0$\\
    Pauli $X$ & $1.41421$ & $2.69\times10^{-27}$ & 2 & $0$\\
    \bottomrule
  \end{tabular}
\end{table}

The corresponding residual-norm upper bounds are approximately
$2.73\times10^{-14}$ and $5.19\times10^{-14}$.  Pauli $X$ succeeds despite its
nonzero charge commutator, directly falsifying charge conservation as a
necessary condition for unconstrained spectral state synthesis.

\subsection{Coupled graph-of-map fixture}

The current coupled experiment fits a complete two-dimensional basis at a
common eigenvalue on one shared geometry.  For a fixed generic $2\times2$
unitary, the default run produced the metrics in~\cref{tab:coupled}.

\begin{table}[htbp]
  \centering
  \caption{Deterministic coupled operator-transfer fixture.}
  \label{tab:coupled}
  \begin{tabular}{lr}
    \toprule
    Metric & Recorded value\\
    \midrule
    full-$W$ squared residual & $7.36\times10^{-17}$\\
    stellar growth steps & $7$\\
    intrinsic boundary drift & $0$\\
    pinned restriction error & $0$\\
    maximum isolated-boundary squared residual & $1.22\times10^{-30}$\\
    recovered-operator relative error & $2.51\times10^{-16}$\\
    maximum held-out full-$W$ squared residual & $1.36\times10^{-16}$\\
    maximum held-out boundary error & $1.84\times10^{-16}$\\
    common-eigenvalue spread & $2.88\times10^{-10}$\\
    \bottomrule
  \end{tabular}
\end{table}

The operator readout and held-out boundary restrictions are at floating-point
precision.  The full-$W$ objective corresponds to a residual-norm bound of
approximately $8.58\times10^{-9}$, so it should not be described as a
machine-precision eigenpair.  The evidence supports accurate transfer in this
fixture, not exact numerical stationarity.

\subsection{Completed stress test}

A completed, confirmed smoke campaign contains 77 cases spanning named gates,
random unitary, normal, non-normal, scaled, singular, charge-preserving,
charge-permuting, and charge-mixing matrices; and a sequence of increasingly
conditioned input frames.  The status counts are in~\cref{tab:stress-status}.

\begin{table}[htbp]
  \centering
  \caption{Outcomes of the completed 77-case stress test.  ``Budget'' means the
  configured optimizer budget ended before certification; it is not a proof of
  nonexistence.}
  \label{tab:stress-status}
  \begin{tabular}{lrr}
    \toprule
    Status & Count & Fraction\\
    \midrule
    verified & 43 & $55.8\%$\\
    optimization budget exhausted & 22 & $28.6\%$\\
    boundary target rejected & 8 & $10.4\%$\\
    operator readout unstable & 3 & $3.9\%$\\
    fitted span not at one common eigenvalue & 1 & $1.3\%$\\
    \midrule
    total & 77 & $100\%$\\
    \bottomrule
  \end{tabular}
\end{table}


Five rejections were zero training outputs, which are valid for singular maps
but incompatible with the current witness normalization.  One input and two
outputs failed the isolated-boundary eigenstate precondition.
No case changed a frozen boundary edge or a pinned boundary restriction.  Among
verified cases, the largest recorded squared residual was
$9.89\times10^{-17}$, the largest relative operator error was
$2.64\times10^{-10}$, the largest held-out relative operator error was
$3.89\times10^{-16}$, the largest held-out full squared residual was
$9.83\times10^{-16}$, and the largest common-eigenvalue spread was
$1.77\times10^{-9}$.  Verified operator condition numbers ranged from $1$ to
$5.45\times10^{16}$ and operator norms from $0.141$ to $10^3$; these ranges do
not imply uniform coverage within the interval.

The largest verified training readout error came from a nearly collinear input
frame.  The conditioning sequence in~\cref{tab:conditioning} shows the expected
loss of identifiability.

\begin{table}[htbp]
  \centering
  \caption{Input-frame conditioning and operator readout.}
  \label{tab:conditioning}
  \begin{tabular}{rrl}
    \toprule
    input condition number & relative readout error & status\\
    \midrule
    $2\times10^2$ & $3.95\times10^{-15}$ & verified\\
    $2\times10^4$ & $2.67\times10^{-12}$ & verified\\
    $2\times10^6$ & $2.64\times10^{-10}$ & verified\\
    $2\times10^8$ & $1.17\times10^{-8}$ & unstable\\
    $2\times10^{10}$ & $9.63\times10^{-7}$ & unstable\\
    $2\times10^{12}$ & $1.88\times10^{-4}$ & unstable\\
    \bottomrule
  \end{tabular}
\end{table}

A separate deterministic set of nine named unitaries verified all nine, with
recorded squared residuals between $6.26\times10^{-17}$ and
$9.11\times10^{-17}$, zero boundary violations, and at most seven growth steps.
Two larger shuffled all-attachment campaigns were still running at the freeze
date of this draft.  No conclusion in this manuscript depends on their partial
records.

These results support three limited conclusions.  First, the coupled solver can
recover one operator from a complete basis and transfer held-out states on the
same frozen geometry.  Second, neither charge conservation nor unitarity is a
necessary numerical precondition.  Third, conditioning and common-eigenvalue
certification are genuine failure modes.  The results do not establish that the
optimizer succeeds for every map or that every successful network is a valid
simplicial cobordism.

\section{Certification and stability}

Let $L$ be Hermitian and let $E_\lambda$ be an eigenspace separated from the
rest of the spectrum by
\begin{equation}
  \delta=\dist(\lambda,\spec(L)\setminus\{\lambda\})>0.
\end{equation}
For any unit vector $v$,
\begin{equation}
  \dist(v,E_\lambda)
  \leq \frac{\norm{(L-\lambda I)v}}{\delta}.
  \label{eq:residual-gap}
\end{equation}
Thus a small eigenresidual has geometric meaning only relative to a spectral
gap.  Subspace versions are instances of Davis--Kahan perturbation bounds
\cite{davis1970}; residual-based Hermitian bounds are discussed, for example,
by Stewart~\cite{stewart1991}.

Operator readout has a separate conditioning factor.  Suppose an exact
boundary frame has blocks $X$ and $Y=UX$, whereas the computed frame has
$\widehat X=X+E$ and $\widehat Y=UX+F$.  Then
\begin{equation}
  \widehat U-U=(F-UE)(X+E)^{-1}.
  \label{eq:readout-error}
\end{equation}
If $\norm{X^{-1}E}<1$, the norm satisfies
\begin{equation}
  \norm{\widehat U-U}
  \leq
  \frac{(\norm F+\norm U\,\norm E)\norm{X^{-1}}}
       {1-\norm{X^{-1}E}}.
  \label{eq:readout-bound}
\end{equation}
This explains why small witness residuals can coexist with an unstable matrix
readout when training inputs are nearly dependent.

A defensible numerical certificate should report at least:

\begin{enumerate}[label=(\arabic*)]
  \item full-complex residual norms and whether stored metrics are squared;
  \item the common-eigenvalue spread across basis witnesses;
  \item the spectral gap separating the intended cluster;
  \item dimension of the boundary trace and singular values of its input
  projection;
  \item boundary geometry drift and pinned restriction error;
  \item relative operator readout error in a declared framing;
  \item held-out linear-combination errors without re-optimizing geometry or
  attachment; and
  \item optimizer budget, restart count, and topology growth.
\end{enumerate}

Permutation enumeration should additionally report both the best and the
distribution over framings.  A best-case result selected after observing the
test input is not an operator certificate.

\section{Composition and transition amplitudes}

Graphs of maps compose as linear relations:
\begin{equation}
  \Gamma_V\circ\Gamma_U
  =\{(x,z):\exists y,\ y=Ux,\ z=Vy\}
  =\Gamma_{VU}.
  \label{eq:relation-composition}
\end{equation}
This identity is the algebraic core of a cobordism composition law.  It does
not yet prove that gluing two triangulated realizations and eliminating the
interface produces the required spectral relation without spurious modes.
That statement requires an additional gluing theorem controlling the interface
Schur complement and eigenvalue multiplicity.

When a framed spectral cobordism $W$ realizes a unique operator, define
\begin{equation}
  \mathfrak Z_\lambda(W)=U
  \quad\text{when}\quad
  \mathcal R_W(\lambda)=\Gamma_U.
  \label{eq:operator-valued-Z}
\end{equation}
For an incoming state $\ket{\psi_B}$ and outgoing test state
$\ket{\psi_A}$, the induced scalar transition amplitude is then
\begin{equation}
  Z(W_{AB};\psi_A,\psi_B)
  =\bra{\psi_A}\mathfrak Z_\lambda(W_{AB})\ket{\psi_B}
  =\matrixel{\psi_A}{U}{\psi_B}.
  \label{eq:amplitude}
\end{equation}
In oriented cobordism notation one may write
\begin{equation}
  W_{AB}=\geo(U),
  \qquad
  \partial W_{AB}
  =\overline{\geo(\psi_B)}\sqcup\geo(\psi_A),
  \label{eq:cobordism-notation}
\end{equation}
where $\sqcup$ denotes disjoint union and the bar reverses the incoming
orientation.  Equations~\eqref{eq:amplitude}--\eqref{eq:cobordism-notation} are
valid as notation for the extracted spectral map.  They are not presently a
derivation of a path integral or a proof that $\mathfrak Z_\lambda$ is a TQFT
functor.  Such a claim would require, at minimum, identity, gluing, monoidality,
and invariance under the chosen equivalence of cobordisms.

\section{Discussion}

The graph-of-map formulation settles the basic semantics of input substitution.
The learned object is neither an isolated boundary eigenvector nor the bulk
with its boundary literally subtracted.  It is the boundary relation obtained
after the bulk has been eliminated.  The bulk acts as a realization of a Schur
response, and the operator is recovered from the kernel of that response.

This perspective also resolves the proposed Choi interpretation.  A Choi vector
stores all matrix entries and can serve as a training target, but promotion to
an operator requires a declared tensor factorization and reshaping convention.
By contrast, $\Gamma_U$ already has operational input and output projections.
The two descriptions are algebraically interconvertible after the frame has
been fixed, but they answer different experimental questions.

The exact response-synthesis theorem is nontrivial and broad: a scalar
positive-weight magnetic network can realize the graph of every finite complex
linear map while preserving arbitrary intrinsic boundary geometry.  It also
shows that charge conservation cannot be the general existence criterion.
Charge instead determines when the relation and its penalty respect a selected
symmetry.  This distinction matters for both theorem statements and numerical
test design.

The scientific limitation is equally clear.  The current universal theorem is
a theorem about finite magnetic graphs, which are discrete geometries but need
not be manifold cobordisms.  The relaxation operates in a narrower class of
triangulated complexes and has finite budgets.  Numerical success across many
operators cannot replace the relative simplicial realization theorem, and
budget exhaustion cannot refute it.  The most valuable next mathematical step
is therefore a response-preserving thickening construction or a proof that
such a construction is obstructed for some boundary data.

\section{Conclusion}

A finite-dimensional operator can be represented intrinsically as the graph of
a boundary spectral relation.  Complete-basis fitting at a common eigenvalue
then supports every new input by linearity, provided the boundary trace has the
correct dimension and a nonsingular input projection.  Every such relation has
a canonical Hermitian penalty, and every Hermitian boundary response admits an
explicit realization by a finite simple scalar magnetic graph with unchanged
intrinsic boundary geometry.  Charge conservation is exactly equivariance of
this relation; attachment permutations select its frame.

The Tessera experiments implement a fixed-boundary inverse spectral relaxation
consistent with these semantics and demonstrate high-accuracy operator readout
and held-out transfer in verified cases.  They also expose optimization,
common-eigenvalue, and conditioning failures.  Establishing universal
realization within nondegenerate triangulated cobordisms, and proving a gluing
law sufficient for a functorial amplitude assignment, remain open.  These
boundaries separate the proven mathematical result from the proposed geometric
generalization.

\appendix

\section{Reproducibility protocol}

The reference implementation is the open-source Tessera repository
\cite{tessera}.  The principal executable examples are:

\begin{itemize}[leftmargin=2.2em]
  \item \path{examples/cobordism/geometric_operators.py}; and
  \item \path{examples/cobordism/geometric_operator_sweep.py}.
\end{itemize}

They are backed by the coupled solver in \texttt{MultiCobordism} and
\texttt{EigenstateSynthesis}.  After building the Python bindings according to
the repository instructions, the deterministic fixture and the completed
77-case record are reproduced by the following invocations:

\begin{verbatim}
python examples/cobordism/geometric_operators.py
python examples/cobordism/geometric_operator_sweep.py \
  --profile smoke --jobs 4 --confirmation-attempts 2 \
  --output /tmp/cobordism/geometric_operator_sweep_smoke_confirmed.json
\end{verbatim}

The larger exhaustive campaign enumerates all six attachment permutations and
uses independently seeded operator generation, schedule shuffling, input
states, and optimizer generation.  It checkpoints every ten cases and retains
every status rather than discarding failures.  Reproduction should record the
repository revision, compiler and linear-algebra library, floating-point
precision, all random seeds, and the complete JSON result.  Confirmatory runs
must preserve the frozen boundary and framing while changing optimizer seeds.

\section{Claim ledger}

\begin{table}[H]
  \centering
  \caption{Epistemic status of the principal claims.}
  \label{tab:claim-ledger}
  \begin{tabular}{>{\raggedright\arraybackslash}p{0.53\textwidth}
                  >{\raggedright\arraybackslash}p{0.20\textwidth}
                  >{\raggedright\arraybackslash}p{0.17\textwidth}}
    \toprule
    Claim & Status & Location\\
    \midrule
    A complete basis at one common eigenvalue transfers every superposition.
      & proved & \cref{thm:basis-span}\\
    $K_U$ has kernel $\Gamma_U$ for every complex linear map.
      & proved & \cref{eq:KU-energy}\\
    Every Hermitian boundary response is realizable by a finite simple magnetic
    graph with fixed induced boundary graph.
      & proved & \cref{thm:response-synthesis}\\
    Charge conservation is invariance of $\Gamma_U$ and commutation of $K_U$
    with boundary charge.
      & proved & \cref{thm:charge}\\
    The current relaxation can recover and transfer multiple operator families.
      & numerical evidence & \cref{sec:numerics}\\
    Every target response admits a valid nondegenerate triangulated-cobordism
    realization with the same frozen boundary.
      & open & \cref{open:relative-simplicial}\\
    Geometric gluing realizes operator composition without spurious modes.
      & open & \cref{eq:relation-composition}\\
    The construction defines a TQFT or path integral.
      & not claimed & \cref{eq:amplitude}\\
    Charge conservation is necessary for unconstrained operator realization.
      & false & \cref{thm:charge} and Table~\ref{tab:historical}\\
    A fitted Choi-like vector alone is an operator accepting new inputs.
      & false without promotion & \cref{eq:operator-readout}\\
    \bottomrule
  \end{tabular}
\end{table}

\begin{thebibliography}{99}

\bibitem{atiyah1988}
M.~F. Atiyah,
\newblock Topological quantum field theories,
\newblock \emph{Publications Math\'ematiques de l'IH\'ES} \textbf{68}
  (1988), 175--186.
\newblock \href{https://doi.org/10.1007/BF02698547}{doi:10.1007/BF02698547}.

\bibitem{jamiolkowski1972}
A.~Jamio\l kowski,
\newblock Linear transformations which preserve trace and positive
  semidefiniteness of operators,
\newblock \emph{Reports on Mathematical Physics} \textbf{3} (1972), 275--278.
\newblock \href{https://doi.org/10.1016/0034-4877(72)90011-0}
  {doi:10.1016/0034-4877(72)90011-0}.

\bibitem{choi1975}
M.-D. Choi,
\newblock Completely positive linear maps on complex matrices,
\newblock \emph{Linear Algebra and its Applications} \textbf{10} (1975),
  285--290.
\newblock \href{https://doi.org/10.1016/0024-3795(75)90075-0}
  {doi:10.1016/0024-3795(75)90075-0}.

\bibitem{dorfler2013}
F.~D\"orfler and F.~Bullo,
\newblock Kron reduction of graphs with applications to electrical networks,
\newblock \emph{IEEE Transactions on Circuits and Systems I} \textbf{60}
  (2013), 150--163.
\newblock \href{https://doi.org/10.1109/TCSI.2012.2215780}
  {doi:10.1109/TCSI.2012.2215780};
  \href{https://arxiv.org/abs/1102.2950}{arXiv:1102.2950}.

\bibitem{friedlander2017}
L.~Friedlander,
\newblock The Dirichlet-to-Neumann operator for quantum graphs,
\newblock preprint (2017),
  \href{https://arxiv.org/abs/1712.08223}{arXiv:1712.08223}.

\bibitem{fallat2024}
S.~M. Fallat, H.~Gupta, and J.~C.-H. Lin,
\newblock The inverse eigenvalue problem for Laplacian matrices of a graph,
\newblock preprint (2024),
  \href{https://arxiv.org/abs/2411.00292}{arXiv:2411.00292}.

\bibitem{fabila2018}
J.~S. Fabila-Carrasco, F.~Lled\'o, and O.~Post,
\newblock Spectral gaps and discrete magnetic Laplacians,
\newblock \emph{Linear Algebra and its Applications} \textbf{547} (2018),
  183--216.
\newblock \href{https://doi.org/10.1016/j.laa.2018.02.006}
  {doi:10.1016/j.laa.2018.02.006};
  \href{https://arxiv.org/abs/1710.01157}{arXiv:1710.01157}.

\bibitem{singer2012}
A.~Singer and H.-T. Wu,
\newblock Vector diffusion maps and the connection Laplacian,
\newblock \emph{Communications on Pure and Applied Mathematics} \textbf{65}
  (2012), 1067--1144.
\newblock \href{https://doi.org/10.1002/cpa.21395}{doi:10.1002/cpa.21395};
  \href{https://arxiv.org/abs/1102.0075}{arXiv:1102.0075}.

\bibitem{pachner1991}
U.~Pachner,
\newblock P.L. homeomorphic manifolds are equivalent by elementary shellings,
\newblock \emph{European Journal of Combinatorics} \textbf{12} (1991),
  129--145.
\newblock \href{https://doi.org/10.1016/S0195-6698(13)80080-7}
  {doi:10.1016/S0195-6698(13)80080-7}.

\bibitem{davis1970}
C.~Davis and W.~M. Kahan,
\newblock The rotation of eigenvectors by a perturbation. III,
\newblock \emph{SIAM Journal on Numerical Analysis} \textbf{7} (1970), 1--46.
\newblock \href{https://doi.org/10.1137/0707001}{doi:10.1137/0707001}.

\bibitem{stewart1991}
G.~W. Stewart,
\newblock Two simple residual bounds for the eigenvalues of a Hermitian matrix,
\newblock \emph{SIAM Journal on Matrix Analysis and Applications} \textbf{12}
  (1991), 205--208.
\newblock \href{https://doi.org/10.1137/0612016}{doi:10.1137/0612016}.

\bibitem{tessera}
A.~Kelleher,
\newblock Tessera: computational discrete differential geometry,
\newblock software repository,
  \href{https://github.com/akellehe/tessera}{github.com/akellehe/tessera}
  (accessed August 26, 2026).

\end{thebibliography}

\end{document}
